\documentclass[a4paper,11pt]{article}
\usepackage[utf8]{inputenc}
\usepackage[T1]{fontenc}
\usepackage[ngerman]{babel}
\usepackage{amsmath,amssymb,amsthm}
\usepackage{xcolor}
\usepackage{tikz}
\usepackage[most]{tcolorbox}
\usepackage{enumitem}
\usepackage{titlesec}
\usepackage[a4paper,margin=2.2cm]{geometry}

\definecolor{bgdark}{HTML}{1E1E1E}
\definecolor{fgtext}{HTML}{E6E6E6}
\definecolor{accent}{HTML}{6CB6FF}
\definecolor{gridcol}{HTML}{4A4A4A}
\definecolor{panel}{HTML}{2A2A2A}
\definecolor{gruen}{HTML}{7BD88F}
\definecolor{gelb}{HTML}{E6C07B}
\definecolor{rot}{HTML}{E06C75}
\pagecolor{bgdark}
\color{fgtext}

\titleformat{\section}{\Large\bfseries\color{accent}}{\thesection}{0.8em}{}
\titleformat{\subsection}{\large\bfseries\color{accent}}{\thesubsection}{0.7em}{}
\setlength{\parindent}{0pt}
\setlength{\parskip}{0.4em}

\newtcolorbox{defbox}[1]{enhanced,breakable,colback=panel,colframe=accent,coltext=fgtext,
  coltitle=bgdark,colbacktitle=accent,fonttitle=\bfseries,title={#1},boxrule=0.6pt,arc=2pt}
\newtcolorbox{satzbox}[1]{enhanced,breakable,colback=panel,colframe=gruen,coltext=fgtext,
  coltitle=bgdark,colbacktitle=gruen,fonttitle=\bfseries,title={#1},boxrule=0.6pt,arc=2pt}
\newtcolorbox{bspbox}[1]{enhanced,breakable,colback=panel,colframe=gridcol,coltext=fgtext,
  coltitle=fgtext,colbacktitle=gridcol,fonttitle=\bfseries,title={#1},boxrule=0.6pt,arc=2pt}
\newtcolorbox{fehlerbox}{enhanced,breakable,colback=panel,colframe=rot,coltext=fgtext,
  coltitle=bgdark,colbacktitle=rot,fonttitle=\bfseries,title={Typische Fehler},boxrule=0.6pt,arc=2pt}
\newtcolorbox{merkbox}{enhanced,breakable,colback=panel,colframe=gelb,coltext=fgtext,
  coltitle=bgdark,colbacktitle=gelb,fonttitle=\bfseries,title={Merkhilfe},boxrule=0.6pt,arc=2pt}

\renewcommand{\contentsname}{\color{accent}Inhaltsverzeichnis}

\begin{document}

\begin{center}
  {\color{accent}\Huge\bfseries Analysis II}\\[0.3em]
  {\Large Lernskript 01 — Gewöhnliche Differenzialgleichungen}\\[0.4em]
  {\color{gridcol}\small Verstehen $\to$ einordnen $\to$ rechnen}
\end{center}
\vspace{0.4em}
{\color{gridcol}\hrule height 1pt}
\vspace{0.8em}

{\color{fgtext}\tableofcontents}
\vspace{0.6em}
{\color{gridcol}\hrule height 1pt}

%====================================================================
\section{Was ist eine Differenzialgleichung?}

\textbf{Intuition.} Eine gewöhnliche Gleichung wie $x^2=4$ sucht eine \emph{Zahl}.
Eine \emph{Differenzialgleichung} (DGL) sucht dagegen eine ganze \emph{Funktion} $y(x)$.
Der Clou: In der Gleichung steht nicht nur $y$, sondern auch ihre Ableitung $y'$. Eine
DGL beschreibt also einen \emph{Zusammenhang zwischen einer Größe und ihrer Änderungs\-rate}.

Das wichtigste Bild des ganzen Kapitels:
\[
  y' = \alpha\, y \quad\Longleftrightarrow\quad
  \text{,,Die Änderung ist proportional zum Bestand.''}
\]
Das ist Wachstum (Zinsen, Bakterien) bzw.\ Zerfall (Radioaktivität), und die Lösung ist
\emph{immer} eine $e$-Funktion. Wer dieses eine Beispiel verstanden hat, hat den roten Faden.

\begin{defbox}{Grundbegriffe}
\begin{itemize}[leftmargin=1.4em,itemsep=2pt]
  \item Eine \textbf{(gewöhnliche) DGL} ist eine Gleichung, die $x$, die gesuchte Funktion
        $y(x)$ und deren Ableitungen $y',y'',\dots$ verknüpft.
  \item Die \textbf{Ordnung} ist die höchste vorkommende Ableitung
        (nur $y'$: 1.\ Ordnung;\ $y''$: 2.\ Ordnung).
  \item Eine \textbf{Lösung} ist eine Funktion $y(x)$, die die Gleichung für alle $x$ erfüllt.
  \item Die \textbf{allgemeine Lösung} enthält eine freie Konstante $c$ (unendlich viele Kurven).
        Ein \textbf{Anfangswert} $y(x_0)=y_0$ legt $c$ fest $\Rightarrow$ \textbf{spezielle Lösung}.
  \item DGL $+$ Anfangswert $=$ \textbf{Anfangswertproblem (AWP)}.
\end{itemize}
\end{defbox}

\begin{bspbox}{Beispiel: eine Lösung erraten und prüfen}
Behauptung: $y(x)=e^{2x}$ löst $y'=2y$. Probe: $y'=2e^{2x}=2\,y$. \checkmark\;
Auch $y=c\,e^{2x}$ löst es für jedes $c$ — das ist die allgemeine Lösung. Mit $y(0)=5$
folgt $c=5$, also die spezielle Lösung $y=5e^{2x}$.
\end{bspbox}

%====================================================================
\section{Einordnung: die drei Fragen}

Bevor du \emph{irgendetwas} rechnest, ordne die DGL ein. Drei Fragen entscheiden über die Methode.

\begin{defbox}{Die drei Einordnungs-Fragen}
\renewcommand{\arraystretch}{1.3}
\begin{tabular}{@{}p{3.1cm}p{4.1cm}p{7.5cm}@{}}
\textbf{(a) Ordnung?} & 1., 2., \dots & höchste Ableitung: $y'\Rightarrow1$,\ $y''\Rightarrow2$ \\
\textbf{(b) linear?} & linear / nichtlinear &
  \textbf{linear}: $y,y'$ nur ,,pur'' (mal Funktionen von $x$).
  \textbf{nichtlinear}: $y^2,\sqrt y,\sin(y),y\cdot y',\dots$ \\
\textbf{(c) homogen?} & homogen / inhomogen &
  \textbf{homogen}: jeder Term enthält $y$ oder $y'$.
  \textbf{inhomogen}: zusätzlicher ,,Störterm'' rein aus $x$ / konstant \\
\end{tabular}
\end{defbox}

\begin{bspbox}{Einordnen im Kopf}
\begin{itemize}[leftmargin=1.4em,itemsep=2pt]
  \item $y'=x^2y$ — 1.\ Ordnung, \textbf{linear}, \textbf{homogen}.
  \item $y'=\sin(y)$ — 1.\ Ordnung, \textbf{nichtlinear} ($y$ im Sinus), homogen.
  \item $y''=y'+x$ — \textbf{2.\ Ordnung}, linear, \textbf{inhomogen} (Störterm $x$).
  \item $y'=xy+x^2$ — 1.\ Ordnung, linear, \textbf{inhomogen} (Störterm $x^2$).
\end{itemize}
\end{bspbox}

\begin{fehlerbox}
\begin{itemize}[leftmargin=1.4em,itemsep=2pt]
  \item ,,Linear'' verwechseln mit ,,$x$ kommt nur linear vor''. Es geht um $y$ und $y'$!
        $y'=x^2y$ ist linear, obwohl $x^2$ auftaucht.
  \item Den konstanten Störterm übersehen: $y'=y+3$ ist \emph{inhomogen} (die $3$ enthält kein $y$).
\end{itemize}
\end{fehlerbox}

%====================================================================
\section{Entscheidungsbaum: welcher Typ -- welche Methode}

\begin{bspbox}{Strategie}
\ttfamily\small
LINEAR? (y, y' nur pur)\\
|\\
+-- JA --> Faktor vor y konstant?\\
\hspace*{2.2em}|\\
\hspace*{2.2em}+-- ja:  y' = a*y\hspace{2.6em}--> TYP A  (e-Funktion direkt)\\
\hspace*{2.2em}+-- nein: y' = a(x)*y\\
\hspace*{2.2em}\hspace*{4.5em}|\\
\hspace*{2.2em}\hspace*{4.5em}+-- homogen          --> TYP B  (y = c*e\^{}A)\\
\hspace*{2.2em}\hspace*{4.5em}+-- inhomogen (+b(x)) --> TYP C  (Loesungsformel)\\
|\\
+-- NEIN (nichtlinear) --> trennbar?  y' = g(x)*h(y)  --> TYP D
\end{bspbox}

Die folgenden vier Abschnitte sind genau die vier Typen.

%====================================================================
\section{Typ A: $y'=\alpha y$ (linear, konstanter Koeffizient)}

\textbf{Intuition.} Änderung proportional zum Bestand $\Rightarrow$ exponentiell.
$\alpha>0$: Wachstum, $\alpha<0$: Zerfall.

\begin{satzbox}{Lösungsformel (auswendig!)}
\[
  y'=\alpha y,\quad y(x_0)=y_0 \qquad\Longrightarrow\qquad
  \boxed{\,y(x)=y_0\,e^{\alpha (x-x_0)}\,}
\]
Hier wird \emph{nicht} integriert — nur $\alpha,\ x_0,\ y_0$ ablesen und einsetzen.
\end{satzbox}

\begin{satzbox}{Halbwerts- und Verdopplungszeit}
\[
  t_{1/2}=\frac{\ln 2}{|\alpha|}\quad(\alpha<0,\ \text{Zerfall}),
  \qquad
  t_{2}=\frac{\ln 2}{\alpha}\quad(\alpha>0,\ \text{Wachstum}).
\]
\end{satzbox}

\begin{bspbox}{Durchgerechnet: Zerfall}
$100\,$g zerfallen in $10\,$h auf $25\,$g. DGL: $y'=\alpha y,\ y(0)=100$, also $y(t)=100e^{\alpha t}$.
\[
  100e^{10\alpha}=25 \;\Rightarrow\; e^{10\alpha}=\tfrac14
  \;\Rightarrow\; 10\alpha=\ln\tfrac14 \;\Rightarrow\; \alpha=\tfrac{\ln(1/4)}{10}\approx-0{,}1386.
\]
Halbwertszeit $t_{1/2}=\dfrac{\ln 2}{0{,}1386}=5\,$h. (Kontrolle: $100\to50\to25$, zwei
Halbierungen in $10\,$h.)
\end{bspbox}

\begin{fehlerbox}
\begin{itemize}[leftmargin=1.4em,itemsep=2pt]
  \item Bei $x_0\neq0$ das $(x-x_0)$ vergessen: $y'=-2y,\ y(1)=5 \Rightarrow y=5e^{-2(x-1)}$,
        \emph{nicht} $5e^{-2x}$.
  \item Vorzeichen von $\alpha$: $\ln(1/4)<0$ — bei Zerfall muss $\alpha$ negativ sein.
\end{itemize}
\end{fehlerbox}

%====================================================================
\section{Typ B: $y'=a(x)\,y$ (linear, homogen, variabler Koeffizient)}

\textbf{Intuition.} Wie Typ A, aber der Proportionalitätsfaktor ist jetzt eine Funktion $a(x)$.

\begin{satzbox}{Verfahren}
\begin{enumerate}[leftmargin=1.6em,itemsep=2pt]
  \item $A(x)=\displaystyle\int a(x)\,dx$\quad (eine Stammfunktion, ohne Konstante).
  \item Allgemeine Lösung: $\boxed{\,y(x)=c\,e^{A(x)}\,}$.
  \item $c$ aus dem Anfangswert bestimmen.
\end{enumerate}
Typ A ist der Spezialfall $a(x)=\alpha=\text{const}\Rightarrow A(x)=\alpha x$.
\end{satzbox}

\begin{bspbox}{Durchgerechnet}
$y'=2x\,y,\ y(0)=3$. \quad $A(x)=\int 2x\,dx=x^2 \Rightarrow y=c\,e^{x^2}$.
Anfangswert: $y(0)=c\,e^0=c=3 \Rightarrow y(x)=3\,e^{x^2}$.
\end{bspbox}

\begin{fehlerbox}
\begin{itemize}[leftmargin=1.4em,itemsep=2pt]
  \item $A(x)$ ist die \emph{Stammfunktion} von $a(x)$, nicht $a(x)$ selbst.
  \item Die Konstante steht \emph{vor} der $e$-Funktion ($c\,e^{A}$), nicht im Exponenten.
\end{itemize}
\end{fehlerbox}

%====================================================================
\section{Typ C: $y'=a(x)\,y+b(x)$ (linear, inhomogen)}

\textbf{Intuition.} Jetzt kommt ein Störterm $b(x)$ dazu (von außen ,,eingespeist''). Die
Gesamtlösung ist homogene Lösung plus eine partikuläre — beides steckt in einer Formel.

\begin{satzbox}{Lösungsformel}
\[
  \boxed{\,y(x)=e^{A(x)}\Big(y_0+\int_{x_0}^{x} b(u)\,e^{-A(u)}\,du\Big)\,},
  \qquad A(x)=\int a(x)\,dx.
\]
\end{satzbox}

\begin{merkbox}
Stur in 5 Schritten abarbeiten:
\textbf{1.} $a(x),b(x)$ ablesen $\to$ \textbf{2.} $A(x)=\int a\,dx$ $\to$
\textbf{3.} $e^{A}$ und $e^{-A}$ hinschreiben $\to$ \textbf{4.} Integral $\int b\,e^{-A}\,du$
lösen $\to$ \textbf{5.} mit $e^{A}$ multiplizieren und $y_0$ einsetzen.
Oft kürzt sich $e^{-A}$ mit $b$ schön weg.
\end{merkbox}

\begin{bspbox}{Durchgerechnet (1)}
$y'=y+e^{x},\ y(0)=0$. \quad $a=1,\ A=x,\ b=e^x$.
\[
  y=e^{x}\Big(0+\int_0^x e^{u}e^{-u}\,du\Big)=e^x\int_0^x 1\,du=e^x\cdot x = x\,e^x.
\]
Probe: $y'=e^x+xe^x=e^x+y$. \checkmark
\end{bspbox}

\begin{bspbox}{Durchgerechnet (2)}
$y'=x^{-1}y+3x^{3},\ y(2)=0$. \quad $a=x^{-1},\ A=\ln x,\ b=3x^3$. \;
Mit $e^{A}=x$ und $e^{-A}=x^{-1}$:
\[
  y=x\Big(0+\int_2^x 3u^3\cdot u^{-1}\,du\Big)=x\int_2^x 3u^2\,du
   =x\big[u^3\big]_2^x=x(x^3-8)=x^4-8x.
\]
Probe: $y(2)=16-16=0$. \checkmark
\end{bspbox}

\begin{fehlerbox}
\begin{itemize}[leftmargin=1.4em,itemsep=2pt]
  \item Im Integral $e^{-A}$ (Minus!) verwenden, außen $e^{+A}$.
  \item Bei bestimmten Grenzen: die \emph{Integrationsvariable} ist $u$, die Grenze ist $x$
        (nicht $x$ als Integrationsvariable benutzen).
  \item $y_0$ steht \emph{innerhalb} der Klammer, wird also auch mit $e^{A}$ multipliziert.
\end{itemize}
\end{fehlerbox}

%====================================================================
\section{Typ D: Trennung der Variablen (nichtlinear, aber trennbar)}

\textbf{Intuition.} Wenn sich die DGL als $y'=g(x)\cdot h(y)$ schreiben lässt, kann man $x$
und $y$ ,,auf getrennte Seiten'' bringen und beide einzeln integrieren.

\begin{satzbox}{Verfahren}
\begin{enumerate}[leftmargin=1.6em,itemsep=2pt]
  \item Schreibe $y'=\dfrac{dy}{dx}=g(x)\,h(y)$ und trenne:
        $\displaystyle \frac{1}{h(y)}\,dy=g(x)\,dx$.
  \item Beide Seiten integrieren (eine Konstante $c$ genügt).
  \item Nach $y$ auflösen (wenn möglich).
  \item Anfangswert $\to$ $c$.
\end{enumerate}
\end{satzbox}

\begin{bspbox}{Durchgerechnet (1)}
$y'=x\,y^2,\ y(0)=1$.
\[
  \int y^{-2}\,dy=\int x\,dx \;\Rightarrow\; -\frac1y=\tfrac12x^2+c.
\]
$y(0)=1\Rightarrow -1=c$. Also $-\tfrac1y=\tfrac12x^2-1 \Rightarrow y=\dfrac{2}{2-x^2}$.
\end{bspbox}

\begin{bspbox}{Durchgerechnet (2)}
$y'=2x\,e^{y},\ y(0)=0$.
\[
  \int e^{-y}\,dy=\int 2x\,dx \;\Rightarrow\; -e^{-y}=x^2+c.
\]
$y(0)=0\Rightarrow -1=c \Rightarrow e^{-y}=1-x^2 \Rightarrow y=-\ln(1-x^2)$.
\end{bspbox}

\begin{fehlerbox}
\begin{itemize}[leftmargin=1.4em,itemsep=2pt]
  \item Nach dem Integrieren das \textbf{Auflösen nach $y$} nicht vergessen.
  \item Beim Teilen durch $h(y)$ kann die Lösung $h(y)=0$ verloren gehen (z.\,B.\ $y\equiv0$).
  \item Vorzeichen bei $\int e^{-y}dy=-e^{-y}$ beachten.
\end{itemize}
\end{fehlerbox}

%====================================================================
\section{Strategie \& Probe -- der rote Faden}

\begin{merkbox}
\textbf{Goldene Reihenfolge bei jeder Aufgabe:}
\;1) Ordnung / linear / homogen bestimmen
\;$\to$ 2) Typ A/B/C/D festlegen
\;$\to$ 3) Rezept abarbeiten
\;$\to$ 4) Anfangswert einsetzen
\;$\to$ 5) Probe.
\end{merkbox}

\textbf{Probe (kostet 30 Sekunden, rettet Punkte):} Setze deine Lösung in die
\emph{ursprüngliche} DGL ein ($y'$ ableiten, mit rechter Seite vergleichen) und prüfe
zusätzlich $y(x_0)=y_0$.

\begin{satzbox}{Spickzettel}
\renewcommand{\arraystretch}{1.5}
\begin{tabular}{@{}p{4.3cm}p{3.5cm}p{6.6cm}@{}}
\textbf{Form} & \textbf{Typ} & \textbf{Lösung} \\[2pt]
$y'=\alpha y$ & A: linear, konst. & $y=y_0\,e^{\alpha(x-x_0)}$ \\
$y'=a(x)\,y$ & B: linear homogen & $y=c\,e^{A(x)},\ A=\int a\,dx$ \\
$y'=a(x)\,y+b(x)$ & C: linear inhom. & $y=e^{A}\big(y_0+\int b\,e^{-A}\,du\big)$ \\
$y'=g(x)\,h(y)$ & D: trennbar & $\int\frac{dy}{h(y)}=\int g(x)\,dx$ \\
\end{tabular}
\end{satzbox}

\vfill
{\color{gridcol}\hrule height 1pt}
\vspace{0.3em}
{\color{gridcol}\small Zum Selbsttesten: siehe \textbf{abfrage\_01\_dgl.pdf}; Lösungen in
\textbf{abfrage\_01\_dgl\_loesung.pdf}.}

\end{document}
